\section{KT-STABLE-PROPAGATION: genuine stable congruences in unbounded weights}
\label{sec:kt-stable-propagation}

We use the compact unitary group, standard full lattice stabilizer,
and class-number-one realization already fixed in KT-U3-DIM and
KT-U3-ENDO-EXACT. Thus, for $a\geq0$, the coefficient representation
$V_a$ has highest weight $(a,0,-a)$ and dimension $(a+1)^3$, and
the algebraic modular-form space is
\[
                     M_a=V_a^{\Gamma},\qquad |\Gamma|=1296.
                                                               \tag{SP.1}
\]
The scalar center acts trivially on $V_a$. In even degrees
$e=a+1\geq2$, the current-run theorem KT-U3-ENDO-EXACT states that
the nonstable part of $M_a$ has dimension $\lfloor a/3\rfloor$;
its complement consists of stable parameters with a single
cuspidal $\mathrm{GL}_3/E$ constituent. Here
$E=\mathbf Q(\sqrt{-3})$. Those spectral and local-level results
are cited, not reconstructed in this section.

The theorem below concerns actual stable automorphic eigensystems
in this full compact space. It does not assert the converse to
source coverage, namely membership in \texttt{Adm3U\_s} with its
additional \texttt{WitA} and \texttt{Unl} predicates.

\begin{lemma}[Finite-group multiplicity growth]
Fix a prime $\ell>3$ and put $G=\operatorname{PGL}_3(\mathbf F_\ell)$.
Let $\overline V_a$ be the reduction of an algebraic integral
lattice in the representation with highest weight $(a,0,-a)$.
For every simple $\overline{\mathbf F}_\ell[G]$-module $S$,
\[
 [\overline V_a:S]
   =\frac{[\overline{\mathbf F}_\ell[G]:S]}{|G|}(a+1)^3+O_{\ell,S}(a+1).
                                                               \tag{SP.2}
\]
Every displayed leading coefficient is strictly positive.
\end{lemma}

\begin{proof}
Fix compatible lifts of prime-to-$\ell$ roots of unity to complex
roots of unity, and use Brauer characters on $\ell$-regular
elements. The reduction has the same formal weight character
as its integral algebraic lattice. If a lift of such an element
to $\mathrm{GL}_3(\mathbf F_\ell)$ has lifted eigenvalues
$x_1,x_2,x_3$, its Brauer character is therefore the Weyl
character
\[
 \beta_a(g)=
 \frac{\det\bigl(x_i^{a+2},x_i,x_i^{-a}\bigr)_{i=1}^3}
      {\det\bigl(x_i^2,x_i,1\bigr)_{i=1}^3}.
                                                               \tag{SP.3}
\]
The expression at coinciding eigenvalues is its continuous
confluent-determinant value. A lift of an $\ell$-regular
projective element can be chosen $\ell$-regular: its order
divides the projective order times $\ell-1$.

If the three eigenvalues are distinct, the denominator in
(SP.3) is fixed and nonzero while each of the six terms in
the numerator has complex absolute value one. Thus the
character is bounded independently of $a$. If exactly two
eigenvalues coincide, cancel their common Vandermonde factor
by differentiating one row. The derivative of $x^{a+2}$ or
$x^{-a}$ contributes at most a linear factor in $a$, while
the resulting denominator is fixed and nonzero. Hence
$\beta_a(g)=O_g(a+1)$ in this case. If all three eigenvalues
coincide, the projective element is the identity and the
central weight zero gives
\[
                      \beta_a(1)=(a+1)^3.
                                                               \tag{SP.4}
\]
There are only finitely many classes, so the constants are
uniform for fixed $\ell$ on all nonidentity $\ell$-regular
elements.

Here are sufficient algebraic details for extracting
multiplicities, without assuming semisimplicity of
$\overline{\mathbf F}_\ell[G]$. The ordinary modular trace
functions of its simple modules are linearly independent:
extend a relation to the group algebra and use a lift of a
diagonal matrix unit in one factor of its split semisimple
quotient. Traces on an element equal traces on its
$\ell$-regular part, since the commuting $\ell$-part acts
unipotently on each eigenspace of the regular part. Thus a
matrix of simple trace values on regular classes has a
nonzero maximal minor over $\overline{\mathbf F}_\ell$.
Lifting that minor entrywise by Brauer characters gives a
nonzero complex determinant: its reduction is the original
nonzero minor. The simple Brauer characters are therefore
linearly independent over $\mathbf C$ as well.

Let $r_S=[\overline{\mathbf F}_\ell[G]:S]$. The regular
module's Brauer character is $|G|$ at the identity and
zero at every other regular element. Indeed its matrices
are permutation matrices with cycles of length the order
of the element, and the roots of unity in a nontrivial
cycle sum to zero. Subtract from $\beta_a$ the regular
character times $(a+1)^3/|G|$. The resulting class function
has every value $O_\ell(a+1)$. Inverting the fixed nonzero
minor just described expresses every simple multiplicity
difference as $O_{\ell,S}(a+1)$, proving (SP.2).
Finally $r_S>0$, since every simple module is a quotient
of the regular module.
\end{proof}

\paragraph{A crude effective version.}
Set
\[
                  N_G=|G|=\ell^3(\ell^2-1)(\ell^3-1).
\]
The preceding proof gives the explicit lower bound
\[
 [\overline V_a:S]\ \geq
 \frac{(a+1)^3}{N_G}
       -N_G^{\,2N_G^2+4}(a+2).
                                                               \tag{SP.2a}
\]
Here are the estimates, to make its effectiveness auditable.
Since the character has central weight zero, rescale an
eigenvalue triple by its first entry. The resulting ratios
have orders dividing the projective element's order, hence
dividing $N_G$. Distinct such roots of unity have distance
at least $4/N_G$. The ordinary or confluent Weyl determinant
therefore gives
$|\beta_a(g)|\leq N_G^3(a+2)$ for every nonidentity regular
element. In the fixed invertible minor of simple Brauer
characters, every entry and every Galois conjugate has
absolute value at most $N_G$. If the minor has size $r$,
then $r\leq N_G$, and its determinant and cofactors have
absolute values at most
$H=N_G^{2N_G}$. Its nonzero determinant is an algebraic
integer in a cyclotomic field of degree at most $N_G$:
the common order of the relevant prime-to-$\ell$ roots
divides the prime-to-$\ell$ part of $N_G$.
The absolute norm is a nonzero integer, so the absolute
value of our chosen determinant is at least
$H^{-(N_G-1)}$. Every inverse entry is at most
$H^{N_G}=N_G^{2N_G^2}$.
Multiplying the residual character-value bound by the at
most $N_G$ entries in a row gives the error in (SP.2a).

In particular,
\[
                  a\geq N_G^{\,N_G^2+3}
        \quad\Longrightarrow\quad
                       [\overline V_a:S]>a/3
                           \quad\hbox{for every }S.
                                                               \tag{SP.2b}
\]
Indeed $a\geq2$, so $a+2\leq2a$; if
$C=N_G^{2N_G^2+4}$, the lower bound is at least
$a(a^2/N_G-2C)$, and the stipulated $a$ gives
$a^2/N_G\geq N_G C$. Since $N_G>3$, this exceeds $a/3$.
This bound is deliberately far larger than a useful
computational threshold.

\begin{lemma}[Propagation through the fixed Hecke action]
Assume $\ell>3$ splits in $E$. Fix a residue eigensystem
$\lambda$ occurring in some $M_{a_0}$ modulo $\ell$ for
all good split-prime Hecke operators, including both
minuscule directions. Then its composition multiplicity
in the reductions of $M_a$ is at least
\[
                            c_\lambda(a+1)^3-C_\lambda(a+1)
                                                               \tag{SP.5}
\]
for constants $c_\lambda>0$ and $C_\lambda$ independent
of $a$.
\end{lemma}

\begin{proof}
At the split place $\ell$, the integral unitary group is
$\mathrm{GL}_3(\mathbf Z_\ell)$. Choose the integral
algebraic lattice $L_a$ in $V_a$. Central weight zero
means its reduction factors through $G$ above. The image
$\overline\Gamma$ has order prime to $\ell$; consequently
the averaging idempotent
\[
       u=|\overline\Gamma|^{-1}
                      \sum_{\gamma\in\overline\Gamma}\gamma
                   \in\overline{\mathbf F}_\ell[G]
                                                               \tag{SP.6}
\]
makes invariants exact and gives
$(L_a^\Gamma)\otimes\overline{\mathbf F}_\ell=u\overline V_a$.

For each rational split prime $q\ne\ell$, choose once and
for all the rational neighbor representatives defining
the unscaled compact Hecke operator in either direction.
Their denominators are supported at $q$, so they are
integral units at $\ell$. The same representatives work
in every coefficient representation. Reduction therefore
gives a fixed element
\[
                h_{q,\pm}=u\left(\sum_j\overline g_{q,\pm,j}\right)u
                          \in u\overline{\mathbf F}_\ell[G]u
                                                               \tag{SP.7}
\]
whose action on $u\overline V_a$ is the reduced Hecke
operator. Let $H$ be the finite-dimensional algebra
generated by all these fixed elements and the identity
$u$. No weight-dependent conjugation of a computed matrix
is being used.

One may realize $u\overline{\mathbf F}_\ell[G]$ as the
weight-zero algebraic modular forms at the auxiliary level
whose subgroup at $\ell$ reduces to the scalar matrices.
The class-number-one decomposition identifies its finite
set with $\overline\Gamma\backslash G$. The operators
(SP.7) are the usual spherical Hecke operators away from
$\ell$ on that space, so they commute. Alternatively,
commutativity is unnecessary for the following exact-functor
argument; on each actual $M_a$ it is already the usual
commuting Hecke action.

The functor $W\mapsto uW$ from finite $G$-modules to
finite $H$-modules is exact. Since $\lambda$ occurs in
$u\overline V_{a_0}$, at least one simple $G$-module $S$
has $[uS:\lambda]>0$. Applying the functor to a composition
series and using additivity of Jordan--H\"older multiplicities
gives
\[
 [u\overline V_a:\lambda]
       =\sum_S[\overline V_a:S][uS:\lambda]
       \geq[\overline V_a:S_0]
                                                               \tag{SP.8}
\]
for one such $S_0$. Equation (SP.2) proves (SP.5).
\end{proof}

\begin{lemma}[Integral eigenflags]
Let $L$ be a free lattice over a complete DVR $O$, and
let a commuting family of integral endomorphisms of $L$
be simultaneously diagonalizable after a finite extension
of its fraction field. Over a sufficiently large finite
coefficient extension, the composition multiplicity of a
specified residue eigencharacter equals the number of
characteristic-zero eigenlines, counted with their vector-space
multiplicities, whose characters reduce to it.
\end{lemma}

\begin{proof}
There are finitely many independent endomorphisms and
eigencharacters, since they lie in $\operatorname{End}_O(L)$;
choose a finite extension splitting them all. Intersect $L$
with a common eigenline. This is a saturated rank-one
sublattice, invariant under every endomorphism, and the
quotient lattice is free. Repeat on the quotient. This
constructs an invariant flag with free rank-one successive
quotients and the characteristic-zero eigencharacters on
them. Reduction is exact at each step because the
quotients are free. Jordan--H\"older multiplicities of the
resulting residue characters give the assertion. The
flag may begin with all eigenlines in any specified
characteristic-zero invariant direct summand.
\end{proof}

\begin{theorem}[KT-STABLE-PROPAGATION]
Fix a split prime $\ell>3$, an even $e_0\geq2$, and a
simultaneous residual Eisenstein eigensystem in $M_{e_0-1}$.
Use the normalization in which the raw compact minuscule
operator is $q$ times the normalized standard Satake
trace, and the opposite direction is $q$ times its dual
trace. Suppose the seed character satisfies, in both
directions and at every split prime $q\ne\ell$,
\[
                 \lambda(h_{q,\pm})
                       =q^{1-e_0}+q+q^{1+e_0}\pmod\ell.
                                                               \tag{SP.9}
\]
Then, for every sufficiently large even $e$ with
$e\equiv e_0\pmod{\ell-1}$, the stable summand of
$M_{e-1}$ contains a characteristic-zero eigenline
whose normalized standard and dual traces satisfy
\[
                \operatorname{tr}_{q,\pm}
                   \equiv q^{-e}+1+q^e\pmod\lambda_e
                                                               \tag{SP.10}
\]
simultaneously at all these primes, for one coefficient
prime $\lambda_e\mid\ell$.
For the normalized attached Galois representation, this
is the global residual semisimplification
\[
                     1\oplus\omega^e\oplus\omega^{-e}
                              \quad\hbox{over }G_E.
                                                               \tag{SP.11}
\]
The explicit sufficient inequality
$e-1\geq N_G^{N_G^2+3}$ is available from (SP.2b).
\end{theorem}

\begin{proof}
Set $a=e-1$. By (SP.5), the multiplicity of the fixed
residue character $\lambda$ in $M_a$ grows at least
cubically in $a$. The entire endoscopic dimension is
$\lfloor a/3\rfloor$ by KT-U3-ENDO-EXACT. The integral
eigenflag lemma applied to the commuting Hecke operators
therefore forces at least one stable characteristic-zero
eigenline reducing to $\lambda$, once $a$ is sufficiently
large. In fact the number of such stable eigenlines,
counting eigenspace multiplicities after splitting
coefficients, is bounded below by (SP.5) minus
$\lfloor a/3\rfloor$. Semisimplicity here is the usual
characteristic-zero automorphic spectral decomposition
of this finite compact modular-form space.

For $e\equiv e_0\pmod{\ell-1}$, the raw Eisenstein targets
in (SP.9) equal $q^{1-e}+q+q^{1+e}$ modulo $\ell$.
Dividing the raw operator eigenvalue by the unit $q$
gives (SP.10). The single residue character is defined
on the algebra generated by \emph{all} these operators,
so the resulting coefficient prime and eigenline work
simultaneously at all of them. They are not selected
separately for different Frobenius primes.

For the Galois attachment use Shin, \emph{Galois representations
arising from some compact Shimura varieties}, Annals of
Mathematics 173 (2011), Theorem~1.2 and its following
compatible-system paragraph, pp.~1647--1648. It applies to
the regular algebraic conjugate-self-dual cuspidal
rank-three constituents in the stable spectrum and supplies
a continuous semisimple system over a finite coefficient
field with the needed local compatibility. Its extra
slight-regularity hypothesis is confined to even rank.
Use the same unitary normalization as in the established
compact source coverage. No characteristic-zero
irreducibility assertion is needed here. The two
directions account for both prime ideals of $E$ over
each split rational prime. Those degree-one prime ideals,
away from any fixed finite exceptional set, have Frobenius
classes dense in $G_E$ for finite quotients: in a normal
closure over $\mathbf Q$, any desired element fixing $E$
is realized by Chebotarev at a rational prime splitting
in $E$, and a prime above it realizes the desired class.
Consequently (SP.10) gives the same residual trace function
as (SP.11) on all of $G_E$. The representations factor
through a finite common residual quotient. Linear
independence of simple modular trace functions, with
dimension $3<\ell$, then identifies their semisimplifications,
as in KT-TRACE-AVOID. This proves (SP.11).
\end{proof}

\begin{corollary}[KT-PROPAGATION-SMALL-SUMMANDS]
Under the hypotheses of KT-STABLE-PROPAGATION, let
$D_a\subset M_a$ be any family of characteristic-zero
Hecke-invariant subspaces along the stipulated progression,
and suppose
\[
                         \dim D_a=o((a+1)^3).
                                                               \tag{SP.11a}
\]
For all sufficiently large $a$ in that progression, a
Hecke-invariant complement of $D_a$ contains a common
eigenline with the same all-prime residual Eisenstein
character. More precisely the number of such lines in
the complement, counted with vector-space multiplicity,
is at least
\[
             c_\lambda(a+1)^3-C_\lambda(a+1)-\dim D_a.
                                                               \tag{SP.11b}
\]
The constants are those of (SP.5). If $D_a$ contains the
entire endoscopic subspace, these surviving lines are stable.
\end{corollary}

\begin{proof}
The characteristic-zero Hecke operators are simultaneously
diagonalizable. Thus $D_a$ is a sum of subspaces of common
eigenspaces and has a Hecke-invariant complement. Choose
an integral eigenflag to count the total dimensions of
the common eigenspaces whose characters reduce to
$\lambda$. Their intersections with $D_a$ have total
dimension at most $\dim D_a$. Subtract that number from
(SP.5), giving (SP.11b). This dimension argument also
allows $D_a$ itself to have been specified only after
extending scalars: a common eigenspace has an algebraic
basis, so if it is not contained in $D_a$, some algebraic
eigenvector lies outside $D_a$. The bound is positive for
all sufficiently large $a$, since $c_\lambda>0$ and
(SP.11a) holds. All-prime normalization and the Galois
conclusion are unchanged. No compatibility between the
subspaces $D_a$ at different weights is assumed.
\end{proof}

\begin{corollary}[KT-STABLE-LARGE-PRIME]
Let
\[
                      \ell=8361790201>2^{31}.
                                                               \tag{SP.12}
\]
For every sufficiently large $t\geq0$, there is a stable
cuspidal compact eigensystem with
\[
              e=18+(\ell-1)t,
\]
whose normalized Galois representation has residual
semisimplification $1\oplus\omega^{18}\oplus\omega^{-18}$
at a prime above $\ell$. Nevertheless, for \emph{every}
$t\geq0$,
\[
                K_{36+2(\ell-1)t}(\mathbf Z)\{\ell\}=0.
                                                               \tag{SP.13}
\]
Thus removing the endoscopic factors and changing the
split Frobenius prime cannot make the level-one compact
stable comparison a unit in these sufficiently large
weights at this fixed prime above $2^{31}$.
\end{corollary}

\begin{proof}
The standard-library certificate
\texttt{large\_prime\_seed\_certificate.py} proves primality
by trial division, verifies $\ell\equiv1\pmod3$, and
computes exactly
\[
 \frac{B_{37,\chi_{-3}}}{37}
       =-\frac{1701990248643001964415638362402}{3},
 \qquad
 \ell\mid1701990248643001964415638362402.
                                                               \tag{SP.14}
\]
It uses the ordinary Bernoulli recurrence through index
$37$ and the exact Bernoulli-polynomial formula for
the character of conductor three. Fretwell--Roberts
Theorem 3.6, with $N=3,M=1,k=37$, gives a primitive
newform $f\in S_{37}(\Gamma_0(3),\chi_{-3})$ satisfying
\[
              a_q(f)\equiv1+\chi_{-3}(q)q^{36}
                            \pmod\lambda
                         \qquad(q\ne3)
                                                               \tag{SP.15}
\]
for one coefficient prime $\lambda\mid\ell$.
The congruence is not CM: at the inert prime two a
level-three CM form would have eigenvalue zero, whereas
the same certificate gives
\[
                    1-2^{36}\equiv6536635074\ne0\pmod\ell.
                                                               \tag{SP.16}
\]

KT-U3-ENDO-INJECT therefore supplies a genuine endoscopic
eigenline in $M_{17}$. Its unitary parameter is
$1\boxplus\operatorname{BC}(\pi_f)$. At each split prime
the raw eigenvalue is $q(1+a_q(f)/q^{18})$; the dual
trace has the same residual value. Equation (SP.15)
gives exactly (SP.9) with $e_0=18$. An integral primitive
vector on this eigenline reduces to the required seed
character, simultaneously for all Hecke operators.
KT-STABLE-PROPAGATION now gives the stable conclusion.

For the rational component, cite the carried complete
certificate K17-BERNOULLI-8361790201 and K17-REFLECTOR:
\[
             B_{\ell-19}\equiv1827838999\ne0\pmod\ell,
             \qquad A_\ell^{[-18]}=0.
\]
Every $e=18+(\ell-1)t$ has the same class-group character
index. The carried comparison (K2), valid in all positive
degrees $K_{4k}$, proves (SP.13). No large-index Bernoulli
enumeration is claimed to have been rerun in this section.
\end{proof}

\paragraph{Limits of the conclusion.}
The theorem gives an explicit, enormous sufficient weight
bound; no smallest weight threshold is computed. It gives actual
stable cuspidal automorphic systems with global residual
congruences, not just an eigenvalue coincidence in an
endoscopic summand. It gives neither a nonzero rational
Selmer class nor automatic membership in the stricter
archive admissible-source predicate. It leaves open the
possibility of eliminating the prime by a different
auxiliary field or by the already successful reflected
class-group test. For this fixed field and the entire
stable compact space, however, an all-weight unit claim
is false even when all split test primes are available.

